\documentclass[10pt,oneside]{report}
\usepackage{soc2style}
\begin{document}
\doctitle{SOC2 Audit and Testing Handbook}{Design assessment, period testing, workpapers and exception handling}{SOC2-07}
\section*{Operating convention and boundaries}
An AICPA criterion identifies an outcome against which controls may be evaluated. The selected point-of-focus headings identify relevant characteristics; they are not discrete certification checklist items. Every owner assignment, technological mechanism and audit test in this suite is a proposed example to be tailored to the system description, commitments, risk profile, period, and service organization responsibilities.

\begin{itemize}
\item Always establish system boundaries, the in-scope trust services categories, service commitments and system requirements first.
\item For a security-only engagement, use the common criteria. For additional trust services categories, include the corresponding additional criteria, subject to properly justified applicability.
\item Document complementary user entity controls and subservice organizations; do not treat vendor assurance reports as replacements for management ownership.
\item Privacy points of focus marked [P] or [C] describe typical processor or controller relevance; assess actual responsibilities rather than treating the annotations as absolute.
\end{itemize}

\chapter{SOC 2 examinations and evidence assertions}
The AICPA source describes a SOC 2 Type 2 examination as reporting on suitability of design and operating effectiveness of controls throughout a specified period, with detailed auditor tests and results. The Type 1 report does not contain an opinion on operating effectiveness or those detailed testing results. This handbook provides management readiness and example auditor-oriented procedures; only an appropriately qualified independent practitioner determines attestation tests and opinions.

\subsection{Four different questions}
\begin{itemize}
\item Criterion applicability: Does this criterion belong to the selected trust services category and scoped system?
\item Design suitability: Is the proposed control likely to meet relevant objectives if implemented as designed?
\item Implementation: Does the described control actually exist in the system?
\item Operating effectiveness: Did it operate as intended throughout the specified period with reliable evidence?
\end{itemize}

\chapter{Testing protocol}
\begin{enumerate}
\item Identify selected trust services categories, reporting period, system boundaries and material commitments.
\item Map criterion to real organization control(s) and identify the individual test assertion for each.
\item Document source population, data owner, extraction controls, exclusions and completeness reconciliation.
\item Select inspection, inquiry, observation, re-performance or corroboration to address the assertion and relevant risks.
\item For automated controls, assess key configurations, change restrictions, dependencies and operational exceptions.
\item Select samples by documented professional judgment or a warranted statistical method; do not claim an AICPA-mandated sample size from this source.
\item Record expected attribute, actual result, timestamp, sample, source link, reviewer, deviation and inquiry resolution.
\item Evaluate exceptions by root cause, frequency, scope, compensating controls, possible pervasive risk and remediation validation.
\item Communicate deficiencies through responsible management and the service auditor as required by engagement arrangements.
\end{enumerate}

\subsection{Illustrative sample considerations (not AICPA sampling policy)}
Sampling depends on the control nature, expected deviation rate, reliability of automated enforcement, population volume, period, inherent risk, control frequency and desired assurance. An automated control may be examined through configuration plus change and exception logs; a monthly manual review may require selecting individual monthly executions. No numeric sample sizes are dictated by the uploaded criteria document.

\chapter{Criterion-specific example tests}
These 61 procedures are proposed only. They must be changed to match the actual control design; testing a non-existent hypothetical control would provide no assurance.

\section{Control Environment}
\Needspace{4\baselineskip}
\subsection*{CC1.1: Ethics and conduct}
Source reference: AICPA TSP 100, p. 14. Management control example: Publish and annually attest a code of conduct; operate an independent reporting and escalation path.

Example procedure: Select hires and incumbent population; inspect communication, acknowledgment and deviations.

Suggested supporting records: Approved conduct policy; attestations; case register with redactions. Potential error pattern: An unsigned policy without communicated expectations does not show operation.

\Needspace{4\baselineskip}
\subsection*{CC1.2: Board oversight}
Source reference: AICPA TSP 100, p. 15. Management control example: Schedule risk and control oversight sessions with documented challenges, decisions and assigned actions.

Example procedure: Trace material control exceptions to board reporting and closure.

Suggested supporting records: Board or committee minutes; risk packs; action tracker. Potential error pattern: Management-only approval is not necessarily independent oversight.

\Needspace{4\baselineskip}
\subsection*{CC1.3: Structures and responsibility}
Source reference: AICPA TSP 100, p. 16. Management control example: Maintain organization chart, RACI, control ownership and vendor oversight responsibility.

Example procedure: Sample high-risk controls for assigned accountable owner and segregation.

Suggested supporting records: Org chart; role descriptions; RACI; delegations. Potential error pattern: Unowned inherited cloud services create gaps across boundaries.

\Needspace{4\baselineskip}
\subsection*{CC1.4: Competence and training}
Source reference: AICPA TSP 100, p. 17. Management control example: Use role profiles, background screening where lawful, onboarding, continuing training and succession plans.

Example procedure: Sample privileged roles for qualification and current training.

Suggested supporting records: Role matrix; training records; screening attestations. Potential error pattern: Completion counts alone do not establish role-appropriate capability.

\Needspace{4\baselineskip}
\subsection*{CC1.5: Accountability and performance}
Source reference: AICPA TSP 100, p. 18. Management control example: Tie control responsibilities to performance review, corrective action and incentives with safe escalation.

Example procedure: Trace missed control duties to accountability and follow-up.

Suggested supporting records: Review rubric; acknowledgments; corrective-action logs. Potential error pattern: Do not use incentives that suppress incident reporting.

\section{Information and Communication}
\Needspace{4\baselineskip}
\subsection*{CC2.1: Quality information and inventory}
Source reference: AICPA TSP 100, p. 19. Management control example: Inventory services, repositories, assets, vendors and data flows with ownership, classification and periodic reconciliation.

Example procedure: Reconcile sampled production resources against inventory and classifications.

Suggested supporting records: CMDB; data-flow diagram; classification register; reconciliation logs. Potential error pattern: Unknown cloud resources and shadow repositories make evidence incomplete.

\Needspace{4\baselineskip}
\subsection*{CC2.2: Internal communications}
Source reference: AICPA TSP 100, p. 21. Management control example: Publish control objectives, incident reporting channels, policy changes and security awareness to relevant personnel.

Example procedure: Sample changes and staff for timely notification and comprehension.

Suggested supporting records: Communications log; trainings; incident channel instructions. Potential error pattern: A training slide without attendance or targeting is weak evidence.

\Needspace{4\baselineskip}
\subsection*{CC2.3: External communications}
Source reference: AICPA TSP 100, p. 22. Management control example: Define accurate customer, regulator, vendor and incident contact channels and commitments.

Example procedure: Select applicable customer commitments and inspect external communication.

Suggested supporting records: Service agreements; trust communications; contact playbook. Potential error pattern: Misleading public descriptions can conflict with implemented controls.

\section{Risk Assessment}
\Needspace{4\baselineskip}
\subsection*{CC3.1: Objectives and risk appetite}
Source reference: AICPA TSP 100, p. 24. Management control example: Record system objectives, commitments, risk tolerances and system boundaries before selecting safeguards.

Example procedure: Trace each material commitment into risk identification and controls.

Suggested supporting records: Risk appetite; system description; SLA register. Potential error pattern: Generic security goals cannot prove customer-specific commitments.

\Needspace{4\baselineskip}
\subsection*{CC3.2: Risk identification and analysis}
Source reference: AICPA TSP 100, p. 27. Management control example: Maintain risk inventory, threat scenarios, likelihood and impact, treatment and periodic re-evaluation.

Example procedure: Sample assets/services and verify risk analysis and owners.

Suggested supporting records: Risk register; workshop results; treatment approvals. Potential error pattern: Risk scoring without linked response or reassessment is incomplete.

\Needspace{4\baselineskip}
\subsection*{CC3.3: Fraud risk}
Source reference: AICPA TSP 100, p. 28. Management control example: Model insider misuse, collusion, credential abuse and financial misrepresentation; design anti-fraud controls.

Example procedure: Trace identified fraud scenarios to mitigations and monitoring.

Suggested supporting records: Fraud scenarios; SoD review; whistleblower handling. Potential error pattern: Do not equate fraud assessment with only malware detection.

\Needspace{4\baselineskip}
\subsection*{CC3.4: Significant change assessment}
Source reference: AICPA TSP 100, p. 29. Management control example: Trigger risk reviews for material architecture, acquisition, regulation and vendor changes.

Example procedure: Sample major changes and assess whether risk review preceded deployment.

Suggested supporting records: Change intake; impact assessments; risk decisions. Potential error pattern: Operational changes may invalidate previously sufficient controls.

\section{Monitoring Activities}
\Needspace{4\baselineskip}
\subsection*{CC4.1: Control monitoring}
Source reference: AICPA TSP 100, p. 30. Management control example: Maintain a monitoring plan mixing continuous telemetry and independent targeted reviews.

Example procedure: Compare monitoring coverage and executed reviews to control inventory.

Suggested supporting records: Control monitoring schedule; dashboards; findings. Potential error pattern: Continuous alerts without ownership and review are not evaluations.

\Needspace{4\baselineskip}
\subsection*{CC4.2: Deficiency communication}
Source reference: AICPA TSP 100, p. 31. Management control example: Grade deficiencies, route them to accountable owners, track deadlines and validate closure.

Example procedure: Sample overdue issues for escalation and independent closure verification.

Suggested supporting records: Findings log; escalation tickets; validation artifacts. Potential error pattern: Closing a ticket without confirming risk mitigation is insufficient.

\section{Control Activities}
\Needspace{4\baselineskip}
\subsection*{CC5.1: Risk-based control activities}
Source reference: AICPA TSP 100, p. 31. Management control example: Choose preventive/detective and manual/automated controls tied to documented risks.

Example procedure: Walk through sampled risks and inspect implemented control activity.

Suggested supporting records: Risk-control matrix; approvals; control narratives. Potential error pattern: Checklist compliance does not replace risk-based design.

\Needspace{4\baselineskip}
\subsection*{CC5.2: Technology general controls}
Source reference: AICPA TSP 100, p. 32. Management control example: Establish technology operations, infrastructure, access, development and maintenance governance.

Example procedure: Test technology dependencies of critical automated controls.

Suggested supporting records: ITGC inventory; platform standards; administrative review. Potential error pattern: Strong application controls can fail on ungoverned supporting infrastructure.

\Needspace{4\baselineskip}
\subsection*{CC5.3: Policy deployment}
Source reference: AICPA TSP 100, p. 33. Management control example: Convert approved policy into assigned, timed operating procedures and exceptions governance.

Example procedure: Select policies and trace to accountable operational executions.

Suggested supporting records: Policies; procedures; schedules; exception register. Potential error pattern: Documents do not establish that activities occurred when required.

\section{Logical and Physical Access Controls}
\Needspace{4\baselineskip}
\subsection*{CC6.1: Logical access architecture}
Source reference: AICPA TSP 100, p. 34. Management control example: Federate identities; protect service identities; enforce MFA, RBAC and secured identity architecture.

Example procedure: Inspect identities and authentication design across in-scope assets.

Suggested supporting records: IAM configuration; architecture; identity inventory. Potential error pattern: Service accounts and machine identities are often overlooked.

\Needspace{4\baselineskip}
\subsection*{CC6.2: Provisioning and deprovisioning}
Source reference: AICPA TSP 100, p. 36. Management control example: Bind user creation, unique identities and timely removal to approved joiner-mover-leaver workflow.

Example procedure: Sample additions and leavers from source population to access outcomes.

Suggested supporting records: Access requests; approval log; HR termination feed. Potential error pattern: Ad hoc accounts that bypass central provisioning break the control.

\Needspace{4\baselineskip}
\subsection*{CC6.3: Least privilege review}
Source reference: AICPA TSP 100, p. 36. Management control example: Use least-privilege roles, approval for privilege elevation, periodic access recertification and revocation.

Example procedure: Trace sampled privileged and ordinary access to approved role requirements.

Suggested supporting records: Access review signoffs; role policy; permissions exports. Potential error pattern: An access export alone is not an approved recertification.

\Needspace{4\baselineskip}
\subsection*{CC6.4: Physical access}
Source reference: AICPA TSP 100, p. 37. Management control example: Restrict facilities and backup-media access with appropriate badges, logs and outsourced control coverage.

Example procedure: Inspect applicable facility access and cloud provider responsibility boundary.

Suggested supporting records: Badge access review; facility records; subservice reports. Potential error pattern: Office badges may not cover data-center physical security.

\Needspace{4\baselineskip}
\subsection*{CC6.5: Secure asset disposal}
Source reference: AICPA TSP 100, p. 37. Management control example: Erase, destroy or make unrecoverable data and software before disposition; retain chain of custody.

Example procedure: Sample decommissioned assets to verify destruction and approval.

Suggested supporting records: Sanitization records; disposal certificates; asset ledger. Potential error pattern: Asset recycling without verifiable sanitization leaves residual exposure.

\Needspace{4\baselineskip}
\subsection*{CC6.6: Perimeter and external threat defenses}
Source reference: AICPA TSP 100, p. 38. Management control example: Protect exposed interfaces with restricted network paths, identity controls and hardening.

Example procedure: Select exposed services and test approved ingress and egress posture.

Suggested supporting records: Firewall rules; segmentation map; WAF/config reviews. Potential error pattern: Perimeter safeguards do not eliminate cloud public resource exposure.

\Needspace{4\baselineskip}
\subsection*{CC6.7: Secure data movement}
Source reference: AICPA TSP 100, p. 38. Management control example: Classify data and restrict export, transmission and removable media; enforce suitable encryption.

Example procedure: Sample sensitive data transfers for authorized recipients and protection.

Suggested supporting records: DLP policy; TLS configurations; export approvals. Potential error pattern: Encrypted transmission does not prove authorized disclosure.

\Needspace{4\baselineskip}
\subsection*{CC6.8: Malicious software prevention}
Source reference: AICPA TSP 100, p. 39. Management control example: Prevent unapproved software and monitor malware on supported endpoints and workloads.

Example procedure: Reconcile endpoint population to protection status and exceptions.

Suggested supporting records: EDR coverage; quarantine events; allowlist policy. Potential error pattern: An installed agent with disabled telemetry is not effective protection.

\section{System Operations}
\Needspace{4\baselineskip}
\subsection*{CC7.1: Vulnerability and drift detection}
Source reference: AICPA TSP 100, p. 40. Management control example: Monitor vulnerability disclosures, baseline drift and asset exposure; prioritize and track remediation.

Example procedure: Sample vulnerable assets and follow discovery, ownership and remediation.

Suggested supporting records: Scan results; configuration diffs; risk acceptance tickets. Potential error pattern: CodeQL alerts alone do not cover host, network and cloud exposures.

\Needspace{4\baselineskip}
\subsection*{CC7.2: Anomaly monitoring}
Source reference: AICPA TSP 100, p. 40. Management control example: Collect and correlate infrastructure, application and identity signals with ownership and defined review.

Example procedure: Select detections and inspect triage and tuning evidence.

Suggested supporting records: Detection rules; SIEM alerts; on-call histories. Potential error pattern: Telemetry collection without analysis is not sufficient monitoring.

\Needspace{4\baselineskip}
\subsection*{CC7.3: Security event evaluation}
Source reference: AICPA TSP 100, p. 41. Management control example: Classify events, assess impact, escalate incidents and preserve investigation records.

Example procedure: Sample events through classification and incident decision records.

Suggested supporting records: Triage cases; severity rubric; incident timeline. Potential error pattern: Alert dismissal without documented investigation undermines event evaluation.

\Needspace{4\baselineskip}
\subsection*{CC7.4: Incident response}
Source reference: AICPA TSP 100, p. 42. Management control example: Practice detection, containment, eradication, coordination, communications and post-incident review.

Example procedure: Walk an incident through roles, approvals, communication and remediation.

Suggested supporting records: IR plan; tabletop results; incident cases. Potential error pattern: A response playbook never exercised may be unproven.

\Needspace{4\baselineskip}
\subsection*{CC7.5: Security incident recovery}
Source reference: AICPA TSP 100, p. 44. Management control example: Define restoration sequencing, integrity checks and post-incident validation and lessons learned.

Example procedure: Trace recovered services to tested restoration and business acceptance.

Suggested supporting records: Recovery playbooks; restore tests; PIR action list. Potential error pattern: Restoring service without validating data integrity can prolong impact.

\section{Change Management}
\Needspace{4\baselineskip}
\subsection*{CC8.1: Controlled change management}
Source reference: AICPA TSP 100, p. 45. Management control example: Use PR reviews, required checks, approvals, release provenance, rollback and emergency change handling.

Example procedure: Sample normal and emergency production changes end-to-end.

Suggested supporting records: Pull requests; checks; releases; change tickets. Potential error pattern: Scan results after merge cannot establish a pre-merge control.

\section{Risk Mitigation}
\Needspace{4\baselineskip}
\subsection*{CC9.1: Business disruption mitigation}
Source reference: AICPA TSP 100, p. 47. Management control example: Define plausible disruption scenarios and proportionate mitigation including alternates and financial options.

Example procedure: Trace business-impact findings to mitigations and assigned owners.

Suggested supporting records: BIA; continuity response plan; mitigation decisions. Potential error pattern: Disaster recovery alone does not address every disruption dependency.

\Needspace{4\baselineskip}
\subsection*{CC9.2: Third-party risk}
Source reference: AICPA TSP 100, p. 48. Management control example: Inventory material vendors and assess security, contracts, shared responsibilities and ongoing performance.

Example procedure: Sample critical providers for due diligence and renewal reassessment.

Suggested supporting records: Vendor register; diligence; contracts; monitoring evidence. Potential error pattern: A supplier's SOC 2 report does not automatically cover your obligations.

\section{Additional Criteria for Availability}
\Needspace{4\baselineskip}
\subsection*{A1.1: Capacity management}
Source reference: AICPA TSP 100, p. 50. Management control example: Track utilization and saturation, forecast load and establish thresholds for scaling and architecture changes.

Example procedure: Recompute capacity headroom and inspect change actions on breaches.

Suggested supporting records: Capacity dashboards; forecast; scaling decisions. Potential error pattern: Availability promises need system-specific capacity assumptions.

\Needspace{4\baselineskip}
\subsection*{A1.2: Environmental and backup protections}
Source reference: AICPA TSP 100, p. 50. Management control example: Implement backup, power/environmental protection, recovery infrastructure, restoration plans and monitoring.

Example procedure: Sample backups for success, protected storage and recovery dependencies.

Suggested supporting records: Backup jobs; retention settings; facility/provider evidence. Potential error pattern: Green backup jobs do not prove recoverability.

\Needspace{4\baselineskip}
\subsection*{A1.3: Recovery testing}
Source reference: AICPA TSP 100, p. 52. Management control example: Periodically exercise restoration procedures, including integrity checks, time objectives and follow-up.

Example procedure: Reperform selected recovery steps and compare outcome with defined objectives.

Suggested supporting records: DR exercise plan; restore logs; issue remediation. Potential error pattern: Test scripts without execution results are not recovery evidence.

\section{Additional Criteria for Confidentiality}
\Needspace{4\baselineskip}
\subsection*{C1.1: Confidential information life cycle}
Source reference: AICPA TSP 100, p. 52. Management control example: Identify confidential datasets, restrict permitted use, retain appropriately and review protection needs.

Example procedure: Sample data classes from creation through storage and access controls.

Suggested supporting records: Classification inventory; data access matrix; retention map. Potential error pattern: Confidentiality is broader than personally identifiable information.

\Needspace{4\baselineskip}
\subsection*{C1.2: Confidential disposal}
Source reference: AICPA TSP 100, p. 52. Management control example: Trigger approved deletion, media destruction and vendor disposition at the end of retention.

Example procedure: Sample expired records and verify timely irreversible disposition.

Suggested supporting records: Retention schedule; deletion jobs; vendor certificates. Potential error pattern: Deleting an application row may leave recoverable replicas and backups.

\section{Additional Criteria for Processing Integrity}
\Needspace{4\baselineskip}
\subsection*{PI1.1: Processing specifications}
Source reference: AICPA TSP 100, p. 53. Management control example: Define data contracts, schemas, field semantics, functional outcomes and user-facing specifications.

Example procedure: Trace specifications to implemented validations and published outputs.

Suggested supporting records: Data dictionary; schema registry; acceptance criteria. Potential error pattern: Undocumented source units and populations cause apparent accuracy errors.

\Needspace{4\baselineskip}
\subsection*{PI1.2: Input completeness and accuracy}
Source reference: AICPA TSP 100, p. 54. Management control example: Validate required fields, types, authorization, duplicates and completeness before processing.

Example procedure: Inject invalid and duplicate data in controlled test and inspect handling.

Suggested supporting records: Validation rules; rejects; ingestion logs. Potential error pattern: Schema validation alone may not detect missing transactions.

\Needspace{4\baselineskip}
\subsection*{PI1.3: Processing correctness}
Source reference: AICPA TSP 100, p. 54. Management control example: Use idempotence, error handling, reconciliation and traceability for transformations.

Example procedure: Trace representative inputs through transformations and exception queues.

Suggested supporting records: Job runs; test results; reconciliation outcomes. Potential error pattern: Successful batch exit codes do not prove complete processing.

\Needspace{4\baselineskip}
\subsection*{PI1.4: Output distribution}
Source reference: AICPA TSP 100, p. 55. Management control example: Validate outputs, authorize recipients, preserve integrity and meet delivery expectations.

Example procedure: Trace approved outputs to recipients, timing and end-to-end completeness.

Suggested supporting records: Delivery manifests; approvals; integrity checks. Potential error pattern: Accurate output sent to the wrong party fails the processing objective.

\Needspace{4\baselineskip}
\subsection*{PI1.5: Processing records retention}
Source reference: AICPA TSP 100, p. 56. Management control example: Retain and protect input, intermediate and output records per system specification.

Example procedure: Sample lifecycle records for accuracy, recoverability and defined retention.

Suggested supporting records: Retention policies; transaction logs; storage tests. Potential error pattern: Deleting intermediates without considering auditability can impair assurance.

\section{Additional Criteria for Privacy: Notice}
\Needspace{4\baselineskip}
\subsection*{P1.1: Privacy notice}
Source reference: AICPA TSP 100, p. 57. Management control example: Maintain accurate notices mapped to actual data processing and communicate material updates.

Example procedure: Compare sampled collection paths with notice and update history.

Suggested supporting records: Privacy notice versions; publication history; data map. Potential error pattern: A generic notice can miss actual third-party collection.

\section{Additional Criteria for Privacy: Choice and Consent}
\Needspace{4\baselineskip}
\subsection*{P2.1: Choice and consent}
Source reference: AICPA TSP 100, p. 58. Management control example: Capture consent or lawful choice decisions with purpose, version, timestamp and withdrawal workflow.

Example procedure: Sample preferences to downstream enforcement and change history.

Suggested supporting records: Consent ledger; UX screens; withdrawal logs. Potential error pattern: A checked box does not prove choice persists across integrations.

\section{Additional Criteria for Privacy: Collection}
\Needspace{4\baselineskip}
\subsection*{P3.1: Purpose-limited collection}
Source reference: AICPA TSP 100, p. 59. Management control example: Map personal information fields to authorized purposes and minimize collection by default.

Example procedure: Compare collected fields with documented purpose and approved scope.

Suggested supporting records: Data inventory; form design; API schemas. Potential error pattern: Telemetry may collect extra identifiers outside primary forms.

\Needspace{4\baselineskip}
\subsection*{P3.2: Explicit consent collection}
Source reference: AICPA TSP 100, p. 60. Management control example: Where explicit consent is required, communicate consequences and preserve affirmative evidence.

Example procedure: Sample transactions requiring explicit consent for recorded decision.

Suggested supporting records: Consent artifacts; notices; audit logs. Potential error pattern: Do not assert every data processing purpose requires consent.

\section{Additional Criteria for Privacy: Use, Retention, and Disposal}
\Needspace{4\baselineskip}
\subsection*{P4.1: Purpose limitation}
Source reference: AICPA TSP 100, p. 60. Management control example: Implement use restrictions, lineage and secondary-use approvals for personal information.

Example procedure: Trace secondary uses to approved purpose and technical enforcement.

Suggested supporting records: Purpose map; access policy; data-use approvals. Potential error pattern: Reusing data for analytics may exceed the stated objective.

\Needspace{4\baselineskip}
\subsection*{P4.2: Personal-information retention}
Source reference: AICPA TSP 100, p. 60. Management control example: Define retention clocks, lawful holds, processing-specific schedules and automated review.

Example procedure: Sample aged records against approved retention and legal holds.

Suggested supporting records: Retention schedule; legal holds; purge reports. Potential error pattern: Infinite default retention undermines stated commitments.

\Needspace{4\baselineskip}
\subsection*{P4.3: Personal-information disposal}
Source reference: AICPA TSP 100, p. 61. Management control example: Implement verifiable deletion across primary stores, replicas and processor handoffs.

Example procedure: Trace selected deletion requests through complete affected systems.

Suggested supporting records: Erasure tickets; deletion evidence; vendor confirmations. Potential error pattern: Backups and third parties often complicate deletion completion.

\section{Additional Criteria for Privacy: Access}
\Needspace{4\baselineskip}
\subsection*{P5.1: Data subject access}
Source reference: AICPA TSP 100, p. 61. Management control example: Authenticate applicable requests, locate personal information and respond within documented parameters.

Example procedure: Sample access requests for verification, completeness and response.

Suggested supporting records: DSAR register; identity check; response record. Potential error pattern: Overcollection during identity verification introduces new privacy risk.

\Needspace{4\baselineskip}
\subsection*{P5.2: Correction and amendment}
Source reference: AICPA TSP 100, p. 62. Management control example: Support correction requests, document decisions and propagate changes as applicable.

Example procedure: Trace completed and denied corrections with reason and propagation.

Suggested supporting records: Correction tickets; updated records; vendor notices. Potential error pattern: A frontend correction may leave stale downstream copies.

\section{Additional Criteria for Privacy: Disclosure and Notification}
\Needspace{4\baselineskip}
\subsection*{P6.1: Authorized disclosure}
Source reference: AICPA TSP 100, p. 63. Management control example: Gate outbound personal information sharing by lawful authorization and documented purpose.

Example procedure: Sample disclosures for authorized basis before transfer.

Suggested supporting records: Data-sharing inventory; agreements; consent reference. Potential error pattern: Consent is criterion wording; local legal and engagement context must be assessed.

\Needspace{4\baselineskip}
\subsection*{P6.2: Disclosure accounting}
Source reference: AICPA TSP 100, p. 63. Management control example: Log approved personal information disclosures with recipient, data scope and date.

Example procedure: Reconcile sampled data exports to complete authorization ledger.

Suggested supporting records: Disclosure log; export records; reconciliation. Potential error pattern: Manual spreadsheets often miss API-based transfers.

\Needspace{4\baselineskip}
\subsection*{P6.3: Unauthorized disclosure recording}
Source reference: AICPA TSP 100, p. 64. Management control example: Record suspected and confirmed unauthorized disclosures with details and disposition.

Example procedure: Trace sampled events to documented records and closure rationale.

Suggested supporting records: Incident register; disclosure timeline; investigation records. Potential error pattern: A detection event without privacy evaluation may be overlooked.

\Needspace{4\baselineskip}
\subsection*{P6.4: Privacy vendor commitments}
Source reference: AICPA TSP 100, p. 64. Management control example: Establish vendor privacy commitments, assess adherence and remediate misuse.

Example procedure: Sample processors for commitment, review and corrective action.

Suggested supporting records: DPA/contract clauses; vendor assessments; remediation log. Potential error pattern: A confidentiality clause alone may omit required privacy obligations.

\Needspace{4\baselineskip}
\subsection*{P6.5: Vendor incident notifications}
Source reference: AICPA TSP 100, p. 65. Management control example: Obtain escalation commitments and integrate third-party notifications into IR procedures.

Example procedure: Examine sampled vendor incidents and reporting timelines.

Suggested supporting records: Notification terms; escalation contacts; vendor tests. Potential error pattern: Untested contractual notice paths can delay regulatory actions.

\Needspace{4\baselineskip}
\subsection*{P6.6: Privacy incident notification}
Source reference: AICPA TSP 100, p. 65. Management control example: Assess recipients and deadlines from applicable commitments and laws; document notification decisions.

Example procedure: Trace incidents to legal analysis and timely notification or rationale.

Suggested supporting records: Breach decision tree; counsel review; notices. Potential error pattern: Do not assume every security event is a notifiable privacy breach.

\Needspace{4\baselineskip}
\subsection*{P6.7: Accounting and disclosure to subject}
Source reference: AICPA TSP 100, p. 65. Management control example: Provide appropriate accounting of held and disclosed personal information using verified records.

Example procedure: Sample requests for completeness, identity verification and response.

Suggested supporting records: DSAR package; inventory; disclosure accounting. Potential error pattern: An incomplete disclosure inventory makes accounting inaccurate.

\section{Additional Criteria for Privacy: Quality}
\Needspace{4\baselineskip}
\subsection*{P7.1: Personal data quality}
Source reference: AICPA TSP 100, p. 66. Management control example: Define validity, currency and completeness rules with correction and synchronization.

Example procedure: Sample key attributes for quality controls and exceptions.

Suggested supporting records: Data-quality dashboard; reconciliations; corrected records. Potential error pattern: Security controls do not necessarily establish data accuracy.

\section{Additional Criteria for Privacy: Monitoring and Enforcement}
\Needspace{4\baselineskip}
\subsection*{P8.1: Privacy inquiry and dispute handling}
Source reference: AICPA TSP 100, p. 66. Management control example: Offer intake, investigation, resolution, appeals and trend monitoring for privacy inquiries.

Example procedure: Select disputes and examine receipt, handling, decision and follow-up.

Suggested supporting records: Case register; SLA results; corrective actions. Potential error pattern: Closing complaints without root-cause review hides systemic failures.

\chapter{Reusable audit worksheets}
\subsection{Control design walkthrough}
\begin{itemize}
\item Criterion ID / control ID / relevant objective / risk statement / service commitment.
\item Control performer and independent reviewer, if required by design.
\item Documented steps and expected behavior, including deviations and normal/emergency paths.
\item Evidence source and how completeness/accuracy are supported.
\item Dependencies on ITGC, third parties, customer controls and configuration management.
\item Design result, reviewer, date, unresolved issue and owner.
\end{itemize}

\subsection{Operating test workpaper}
\begin{itemize}
\item Population name, query, totals, dates, scope, reconciliation and preparer/reviewer.
\item Sampling basis, unit identifiers, expected criteria, procedures and supporting source records.
\item Test outcomes for each attribute; deviations, compensating evidence, inquiry, conclusion and follow-up.
\item Link to source export and integrity metadata; redact data as appropriate, preserving auditor access under engagement terms.
\end{itemize}

\subsection{Exception versus deficiency}
An observed control deviation is not automatically a report qualification. Characterize the facts, cause, period, affected objectives, prevalence, alternate controls and management remediation before forming conclusions. Practitioners determine attestation implications; management is responsible for accurate representations and remediation.

\chapter{Independent review and reporting handoff}
\begin{enumerate}
\item Confirm narrative scope, including vendor roles and service commitments, against operational evidence.
\item Reconcile the full population of IT controls and selected trust services criteria to testing workpapers.
\item Separate control failures from evidence-availability issues without minimizing either.
\item Ensure all unresolved exceptions are in a reconciled management issue list.
\item Review whether any assertion depends on incomplete subservice reports or complementary user entity controls.
\item Keep test artifacts access-controlled, retrievable, dated and associated with specific report periods.
\item Escalate legal, financial, regulatory or material assurance questions to counsel, management or licensed practitioner.
\end{enumerate}


\chapter*{Source and editorial protocol}
\addcontentsline{toc}{chapter}{Source and editorial protocol}
This independent document is grounded in the user-supplied \emph{AICPA TSP Section 100: 2017 Trust Services Criteria for Security, Availability, Processing Integrity, Confidentiality, and Privacy (With Revised Points of Focus---2022)}, 75-page PDF, especially paragraphs .01--.29, the criterion table on pp.~14--67, and the glossary. Criterion wording is quoted when labeled as such. Selected point-of-focus labels are aids to interpretation and may be incomplete. Engineering recommendations, role assignments, tools, tests, dates, scorecards, and example frequencies are author-developed examples and are not AICPA requirements. The 2022 revisions changed points of focus, not the 2017 criteria (source p.~3). The source states that an assessment of each point of focus is not required (p.~5). This document does not constitute a CPA attestation or legal advice.
\end{document}
